\section{A Motivating Example}\label{sec:b-example}

\begin{figure}[t]
  \centering\scriptsize
  \renewcommand{\arraystretch}{0.82}
  \begin{tabular}{@{}r@{\hspace{5pt}}l@{\hspace{6pt}}l@{}}
    & \textbf{Program} & \textbf{Slice of \texttt{q[4]}} \\[1pt]
    \color{black!50}3  & \texttt{hops = 4} & C \\
    \color{black!50}4  & \texttt{q = QuantumRegister(1 + 2 * hops, "q")} & C \\
    \color{black!50}5  & \texttt{c = ClassicalRegister(2 * hops, "c")} & C \\
    \color{black!50}6  & \texttt{out = ClassicalRegister(1, "out")} & C \\
    \color{black!50}7  & \texttt{qc = QuantumCircuit(q, c, out)} & C \\
    \color{black!50}8  & \texttt{qc.u(0.4, 0.2, 0.1, q[0])} & Q \\
    \color{black!50}9  & \texttt{src = 0} & C \\
    \color{black!50}10 & \texttt{for h in range(hops):} & C \\
    \color{black!50}11 & \texttt{\ \ a, b = 1 + 2 * h, 2 + 2 * h} & C \\
    \color{black!50}12 & \texttt{\ \ qc.h(q[a])} & Q, $h = 0, 1, 2$ \\
    \color{black!50}13 & \texttt{\ \ qc.cx(q[a], q[b])} & Q, $h = 0, 1, 2$ \\
    \color{black!50}14 & \texttt{\ \ qc.cx(q[src], q[a])} & Q, $h = 0, 1, 2$ \\
    \color{black!50}15 & \texttt{\ \ qc.h(q[src])} & Q, $h = 0, 1, 2$ \\
    \color{black!50}16 & \texttt{\ \ qc.measure(q[src], c[2 * h])} & Q, $h = 0, 1, 2$ \\
    \color{black!50}17 & \texttt{\ \ qc.measure(q[a], c[2 * h + 1])} & Q, $h = 0, 1$ \\
    \color{black!50}18 & \texttt{\ \ with qc.if\_test((c[2 * h + 1], 1)):} & C \\
    \color{black!50}19 & \texttt{\ \ \ \ qc.x(q[b])} & Q, $h = 0, 1$ \\
    \color{black!50}20 & \texttt{\ \ with qc.if\_test((c[2 * h], 1)):} & C \\
    \color{black!50}21 & \texttt{\ \ \ \ qc.z(q[b])} & Q, $h = 0, 1$ \\
    \color{black!50}22 & \texttt{\ \ src = b} & C \\
    \color{black!50}23 & \texttt{qc.measure(q[src], out[0])} & \\
  \end{tabular}
  \caption{A chain of four teleportation hops, our own program (the import on lines 1 and 2 is omitted). The column
  gives the slice of \texttt{q[4]} that \tool{} computes without running the program: Q marks a line whose
  operations are in the slice, with the loop iterations that put it there, and C a statement that computes them.
  srcSlice's profile of \texttt{qc} instead lists every gate line, 8, 12 to 17, 19, 21 and 23.}
  \label{fig:b-example}
\end{figure}

Fig.~\ref{fig:b-example} shows a Qiskit program that teleports the state of \texttt{q[0]} along a chain of qubits. Each
hop $h$ prepares an entangled pair \texttt{q[a]}, \texttt{q[b]} (lines 12 and 13), measures the current carrier
\texttt{q[src]} and \texttt{q[a]} (lines 14 to 17), and repairs \texttt{q[b]} with gates conditioned on the two measured
bits (lines 18 to 21). The carrier then moves to \texttt{q[b]}, so the variable \texttt{src} names a different qubit on
every iteration, and the qubit indices on lines 11 and 22 are computed from the loop counter.

Suppose the state after the second hop, carried by \texttt{q[4]}, looks wrong. A developer who reads the program has to
work out which iterations act on \texttt{q[4]}, follow \texttt{src} through two reassignments, and remember that a
measured bit steers the repair on lines 19 and 21. A classical slicer does not help. srcSlice reads \texttt{qc} as a
single variable that every gate line uses, so its slice of \texttt{qc} is every gate line, and nothing in it
distinguishes one hop from another or one qubit from another.

The slice of \texttt{q[4]} that \tool{} computes keeps \numExSliceOps{} of the \numExAllOps{} operations. It
keeps the first two hops in full and, of the third, the gates on lines 12 to 16, and it names these iterations. The
third hop's preparation of \texttt{q[5]} and \texttt{q[6]} is kept although neither qubit is \texttt{q[4]}, because the
CNOT of line 14 in that hop has \texttt{q[4]} as its control and can change it. The slice also keeps the statements
that compute the indices and the carrier, lines 9 to 11 and 22, and leaves out line 23, which measures another qubit.
None of this needs the program to run, so it works on programs that cannot be run, such as fragments or programs for
an older Qiskit release.
